\chapter{Conclusion and Future Work}
\label{ch:conclusion}

\section{Summary}
This thesis presented \texttt{scrapktm}, a reproducible pipeline for semantic retrieval over Nepali municipal legal documents. The system scrapes and curates a 50-document gold-standard corpus (830 pages, 446 chunks), routes PDFs through a hybrid PyMuPDF/Surya OCR extractor, applies semantic chunking, and benchmarks dense retrieval with FAISS.

The 11-query research bank exposes realistic failure modes---list extraction, table lookup, cross-references, and code-switching---that monolingual English baselines struggle to address. A ground-truth mapping utility diagnoses chunk-ID schema mismatches that otherwise mask retrieval quality entirely.

\section{Key Findings}
\begin{enumerate}[leftmargin=*]
  \item Strict chunk-ID evaluation can produce zero metrics even when retrieval is partially correct.
  \item Kathmandu financial documents yield the highest ground-truth mapping confidence due to clean text and preserved section numbering.
  \item Lalitpur business registration documents suffer from OCR noise and coarse chunking, depressing both mapping and projected recall.
  \item Multilingual embeddings (BGE-M3) with offline caching are essential for Nepali legal RAG on local hardware.
\end{enumerate}

\section{Future Work}
\begin{itemize}[leftmargin=*]
  \item \textbf{Expand corpus and query bank}: Scale to 200+ documents and 50+ annotated queries with inter-annotator agreement.
  \item \textbf{Hybrid retrieval}: Combine dense embeddings with BM25 over Devanagari-tokenized text.
  \item \textbf{Cross-encoder reranking}: Add a second-stage reranker for top-20 candidates.
  \item \textbf{End-to-end RAG evaluation}: Measure answer correctness with LLM-as-judge and human legal review.
  \item \textbf{OCR improvement}: Fine-tune or post-process Surya output for Nepali legal typography.
  \item \textbf{Production deployment}: Package as a municipal legal assistant with citation-backed responses.
  \item \textbf{Unified chunk schema}: Enforce a single chunk-ID convention across annotation, indexing, and UI citation.
\end{itemize}

\section{Reproducibility Statement}
All code, corpus manifests, query banks, and benchmark scripts are available in the \texttt{scrapktm} repository. Researchers can reproduce results by:
\begin{enumerate}[leftmargin=*]
  \item Installing dependencies from \texttt{requirements.txt}.
  \item Running \texttt{extract\_corpus.py} against the curated manifest.
  \item Executing \texttt{ground\_truth\_mapper.py} to align annotations.
  \item Running \texttt{embedding\_benchmark.py} with locally cached models.
\end{enumerate}

The pipeline requires no external API services for retrieval evaluation, supporting fully offline research on Apple Silicon hardware.
